%% Documentation of the BDD-based SAT encodings of fd-fts-sat, written with the JAIR
%% author kit (ACM-based jair.cls, 2025 style). Build:
%%   pdflatex bdd-sat && biber bdd-sat && pdflatex bdd-sat && pdflatex bdd-sat
%% The option tables (options-*.tex) are generated from the C++ sources:
%%   python3 extract_options.py
\documentclass[manuscript, screen]{jair}

\setcopyright{cc}
\acmDOI{}
\JAIRAE{}
\JAIRTrack{}
\acmVolume{}
\acmArticle{}
\acmMonth{10}
\acmYear{2026}

\RequirePackage[
  datamodel=acmdatamodel,
  style=acmauthoryear,
  backend=biber,
  giveninits=true,
  uniquename=init
  ]{biblatex}
\renewcommand*{\bibopenbracket}{(}
\renewcommand*{\bibclosebracket}{)}
\addbibresource{bdd-sat.bib}

\usepackage{booktabs}
\usepackage{longtable}
\usepackage{amsmath}
\usepackage{xspace}

\newcommand{\opt}[1]{\texttt{#1}}
\newcommand{\bddsat}{\opt{bdd\_sat}\xspace}
\newcommand{\B}{\mathcal{B}}
\newcommand{\prev}[1]{\mathit{prev}_{#1}}
\newcommand{\nxt}[1]{\mathit{next}_{#1}}

\begin{document}

\title[BDD-Based SAT Encodings for Factored Transition Systems]{BDD-Based SAT Encodings for
Planning on Factored Transition Systems: Encodings, Label Orders, and Constraints}

% TODO: check the author list and affiliations
\author{Gregor Behnke}
\email{g.behnke@uva.nl}
\affiliation{%
  \institution{University of Amsterdam}
  \city{Amsterdam}
  \country{The Netherlands}}

\renewcommand{\shortauthors}{Behnke}

\begin{abstract}
  This document describes the \bddsat encoder of the planner fd-fts-sat and everything that can be
  configured around it. The planner does not plan on the finite-domain representation of a task but
  on a factored transition system (FTS), possibly after a merge-and-shrink style transformation, and
  encodes bounded plan existence into SAT. For every factor, \bddsat represents the label sets that
  move the factor from one state to another within one time step as binary decision diagrams (BDDs)
  over one variable per label, and translates these BDDs into clauses. We describe how the BDDs are
  built, how they can be strengthened, the alternative CNF translations (Tseitin variants, cube and
  path encodings), the redundant clauses that can be derived from them (label implications, state
  units, leave and enter clauses, and a goal condition on the last step), the label orders that
  define the step semantics, and the search engines and length strategies that use the encoding.
  Every option is listed with its default in the appendix; the tables there are generated from the
  source code.
\end{abstract}

\keywords{Planning as satisfiability, factored transition systems, binary decision diagrams, CNF
encodings, merge-and-shrink}

\maketitle

%% ===========================================================================
\section{Introduction}

Planning as satisfiability \parencite{kautz-selman-ecai1992} answers the question whether a plan of
a given length exists by encoding it as a propositional formula. The planner described here, a fork
of Fast Downward \parencite{helmert-jair2006}, does so for a \emph{factored transition system}: the
finite-domain task is turned into an atomic FTS with one factor per variable, which can be
transformed further by label reduction, shrinking and merging
\parencite{sievers-helmert-jair2021,torralba-sievers-ijcai2019}. Plans for the transformed FTS are
mapped back to the original task afterwards.

The SAT encodings work on the FTS directly. A plan step may contain several labels, applied in a
fixed \emph{label order}; the order thus defines the step semantics. The \bddsat encoder describes,
for every factor and every pair of a source and a target state, the label sets that lead from the
source to the target within one step by a BDD \parencite{bryant-ieee1986}, and translates these BDDs
into clauses. Everything else in this document is a variation of this idea: how the BDDs are built
and strengthened, how they are translated, which consequences of them are added as redundant
clauses, and which label order is used.

The document is organised along the pipeline. Section~\ref{sec:prelim} fixes the notation and the
step semantics. Section~\ref{sec:bdds} describes the construction of the BDDs,
Section~\ref{sec:cnf} their translation into CNF, Section~\ref{sec:cut} the strengthening of the BDDs
by other factors, and Section~\ref{sec:redundant} the redundant clauses. Section~\ref{sec:orders}
describes the label orders, Section~\ref{sec:engines} the search engines and length strategies, and
Section~\ref{sec:empirical} summarises what we observed empirically. The appendix lists every
option.

%% ===========================================================================
\section{Preliminaries}\label{sec:prelim}

A \emph{transition system} is a tuple $\Theta = (S, L, T, s_0, S_\star)$ with a finite set of states
$S$, a finite set of labels $L$, transitions $T \subseteq S \times L \times S$, an initial state
$s_0$ and goal states $S_\star \subseteq S$. A \emph{factored transition system} is a set of
transition systems $\Theta_1, \dots, \Theta_n$ over the same labels. A label is applicable in a
global state $(s_1,\dots,s_n)$ if every factor has a transition with that label from its current
state, and applying it moves every factor along such a transition. A label $\ell$ is
\emph{irrelevant} for a factor if it has exactly the self-loops $(s,\ell,s)$ for all $s$; such labels
do not occur in that factor's BDDs. Note that a label with a self-loop at every state may still move
the factor (it can have further transitions); only labels with nothing but self-loops are
irrelevant.

\paragraph{Step semantics.} A \emph{label order} is a total order $\prec$ on $L$. One time step of a
plan is a set of labels $\{\ell_1 \prec \dots \prec \ell_k\}$, applied as the sequence
$\ell_1,\dots,\ell_k$; each label occurs at most once per step and the empty step is always allowed.
A plan with $m$ actions thus fits into $1 + b$ steps, where $b$ is the number of consecutive pairs
that are not increasing in $\prec$. The label order is therefore a central parameter: it decides how
many time steps a plan needs (Section~\ref{sec:orders}).

\paragraph{Variables.} For every time step $t$ and factor $i$ there is one state variable per state,
with an exactly-one constraint, and for every step one variable per label. We write $\prev{s}$ and
$\nxt{s'}$ for the state variables of a factor at the beginning and end of the step. The initial
state is asserted at the first step and, for every factor, one of its goal states at the last.

%% ===========================================================================
\section{Building the BDDs}\label{sec:bdds}

For every factor $i$ and every pair of states $s, s'$ of that factor, $\B_i(s,s')$ is a BDD over the
label variables that is true exactly for the label sets that move the factor from $s$ to $s'$ in one
step under the step semantics.

\paragraph{Full reachability.} With $\B_i(s,s')$ initialised to $[s = s']$ and the labels processed in
reverse label order, one dynamic-programming step for label $\ell$ is
\[
  \B_i(s,s') \leftarrow \neg\ell \wedge \B_i(s,s') \;\vee\; \bigvee_{(s,\ell,u) \in T_i} \ell \wedge \B_i(u,s'),
\]
i.e., either the label is not used, or it is used first and the remaining (later) labels lead from
its target to $s'$. Afterwards $\B_i(s,s')$ describes arbitrary label sequences that respect the order
(\opt{one\_step\_only=false}). This allows a factor to make several transitions within one step.

\paragraph{One step.} With \opt{one\_step\_only=true}, a factor makes at most one real transition per
step, framed by self-loops: labels before it in the order must self-loop on the source, labels after
it on the target. This is closer to $\exists$-step semantics, gives smaller BDDs for some factors and
much larger ones for others. \opt{bdd\_size\_limit} interpolates: oversized full-reachability BDDs of
a factor are replaced by their one-step counterparts until a node budget is met.

\paragraph{Variable order.} The BDD variable of a label is, by default, its position in the label
order. \opt{bdd\_variable\_order} chooses the BDD variable order independently (any label order
finder), and \opt{bdd\_reorder} lets CUDD \parencite{somenzi-cudd} reorder the variables, once after
construction or dynamically (\opt{bdd\_reorder\_dynamic}). The label order remains the semantics; the
variable order only affects the size of the BDDs. \opt{alternate\_label\_orders} uses two label orders
on alternating steps.

\paragraph{Split and combined BDDs.} By default the encoding keeps one BDD per $(i,s,s')$, conditioned
on the state variables (\emph{split}). With \opt{combinebdds=true} each factor has a single BDD over
the current-state variables $x_0,\dots,x_{n-1}$, the next-state variables $y_0,\dots,y_{n-1}$ and the
labels, which must equal $\B_i(s,s')$ whenever exactly $x_s$ and $y_{s'}$ are true. The original
construction is a disjunction of spelled-out minterms,
\[
  \bigvee_{s,s'} \Bigl( x_s \wedge \bigwedge_{t \neq s} \neg x_t \wedge y_{s'} \wedge \bigwedge_{t' \neq s'} \neg y_{t'} \wedge \B_i(s,s') \Bigr),
\]
which costs $\Theta(|S|^3)$ BDD operations and fixes the value on impossible assignments.
\opt{combine\_ite=true} builds instead the if-then-else chain
\[
  F = \mathit{ite}(x_0, G_0, \mathit{ite}(x_1, G_1, \dots \mathit{ite}(x_{n-1}, G_{n-1}, \bot)\dots)),
  \qquad
  G_s = \mathit{ite}(y_0, \B_i(s,0), \mathit{ite}(y_1, \B_i(s,1), \dots, \bot)\dots),
\]
in which the first true state variable decides. On one-hot assignments, which the exactly-one
constraints guarantee, $F$ coincides with the disjunction above; elsewhere it may differ, which is
sound for the same reason. Because the state variables lie above the labels in the variable order,
every $\mathit{ite}$ only creates a node: $\Theta(|S|^2)$ constant-time steps.

\paragraph{Budgets.} \opt{bdd\_init\_time\_limit} and \opt{bdd\_node\_limit} abort the construction
when exceeded; the CUDD manager is sized by \opt{cudd\_init\_nodes}, \opt{cudd\_cache\_size} and
\opt{cudd\_max\_memory\_mb}.

%% ===========================================================================
\section{Translating the BDDs into CNF}\label{sec:cnf}

In the split representation, the BDD of $(i,s,s')$ is translated with the conditions
$\prev{s} \wedge \nxt{s'}$: its root must be true when both hold. Pairs whose BDD is constant false
yield the clause $\neg\prev{s} \vee \neg\nxt{s'}$. Since every pair, including the impossible ones,
is constrained, no frame axioms are needed. In the combined representation the factor's single BDD is
asserted at every step.

\paragraph{Tseitin.} Every BDD node gets a variable that stands for the sub-BDD below it
\parencite{tseitin1968,een-sorensson-jsat2006}. With \opt{impltseitsin=true} (default) only the
implication from the node variable to its children is encoded, which suffices because the roots are
asserted; \opt{impltseitsin=false} encodes both directions. \opt{omitforcedvariables} skips node
variables for nodes with in-degree at most \opt{forcedvariablesthreshold}, encoding them inline. The
Tseitin variables are shared within a factor and time step, which is valid because they represent the
sub-BDD itself.

\paragraph{Consistency versus domain consistency.} \textcite{abio-et-al-cpaior2016} show that this
translation is only \emph{consistent}: unit propagation detects a violated BDD once all variables are
assigned, but does not derive the values a partial assignment forces. A BDD that implies
$\ell_1 \rightarrow \neg\ell_2$ does not propagate $\neg\ell_2$ when $\ell_1$ is set. This observation
motivates the path encoding and the redundant clauses of Section~\ref{sec:redundant}.

\paragraph{Cube encoding.} With \opt{cube\_encoding\_max=k}, a BDD with at most $k$ paths to the
false terminal is encoded as one clause per such path (its negation), without auxiliary variables.
This is exact and propagation-complete per BDD, but exponential in general.

\paragraph{Path encoding.} \opt{path\_encoding=true} uses NNFPath \parencite{abio-et-al-cpaior2016}.
Variables mean ``this node (edge) lies on the path from the root to the true terminal''. A node on the
path takes one of its outgoing edges; an edge on the path implies its label literal and its child;
node and literal imply the edge; every non-root node on the path needs an incoming edge on the path;
and, for every variable $x$ of the BDD and value $v$,
\[
  (x = v) \;\rightarrow\; \bigvee \{\text{edges carrying } x = v \text{ at } x\text{'s level}\} \vee
  \bigvee \{\text{edges jumping over } x\text{'s level}\}. \tag{P3}
\]
P3 makes unit propagation domain consistent: once no path through $x = v$ is left, $x \neq v$ is
propagated. Three points go beyond the paper:
\begin{itemize}
  \item The paper treats quasi-reduced diagrams. Our BDDs are reduced, so edges may skip levels; P3
        includes such edges for every level they skip. Equivalently one could quasi-reduce the BDD.
  \item CUDD's complement edges are handled by giving every node a parity, so that one CUDD node may
        appear twice in the encoding.
  \item The path variables cannot be shared between diagrams, but CUDD shares nodes heavily between
        the $(s,s')$ BDDs of a factor. \opt{path\_sharing=true} therefore encodes all of a factor's BDDs
        as one diagram with several roots. The root of pair $(s,s')$ is activated by an auxiliary
        variable equivalent to $\prev{s} \wedge \nxt{s'}$; since exactly one pair is active, the
        active nodes still form a single path, and roots below a level count as skipping it in P3.
\end{itemize}
The path encoding uses about three variables per node instead of one.

%% ===========================================================================
\section{Strengthening the BDDs: Cutting}\label{sec:cut}

A label set that factor $j$ cannot perform from any state can never occur in a plan, whatever factor
$i$ allows. With \opt{cutbdds=true}, the union over all pairs of factor $j$'s BDDs, projected onto the
labels relevant to factor $i$ (all other labels existentially quantified), is conjoined into factor
$i$'s BDDs, up to a fixpoint. \opt{cut\_mode} selects how:
\begin{description}
  \item[\opt{AND}] conjunction;
  \item[\opt{AND\_LIMITED}] conjunction with a node limit of \opt{cut\_limit\_factor} times the target's
        size, kept only if the result is not larger;
  \item[\opt{AND\_CAPPED}] the same limit, but the result is kept even if larger;
  \item[\opt{RESTRICT}] \texttt{Cudd\_bddRestrict} with the projection as care set. This is only sound
        when care sets come from unrestricted BDDs, so restriction flows along the factor order only;
        restricting in both directions lets two factors justify each other and accept label sets
        neither allows;
  \item[\opt{GLOBAL\_AND}, \opt{GLOBAL\_AND\_LIMITED}, \opt{GLOBAL\_RESTRICT}] the conjunction $G$ of
        all factors' unions, projected per factor (for restriction, the prefix of earlier factors).
\end{description}
\opt{cut\_min\_overlap} restricts pairwise cutting to factor pairs that share at least the given
fraction of the target's labels, and \opt{coverbdds} reports state-to-label implications for
diagnosis only.

%% ===========================================================================
\section{Redundant Clauses Derived from the BDDs}\label{sec:redundant}

The clauses in this section are logical consequences of the BDDs. They do not change the set of
models; they add propagation the Tseitin translation lacks. All of them are optional and can be
combined.

\subsection{Label implications}

Let $U_i = \bigvee_{s,s'} \B_i(s,s')$ be the label sets factor $i$ permits in one step, from any
state. Every implication over label literals that holds in $U_i$ holds at every step. The essential
variables of a BDD (CUDD's \texttt{FindEssential}) are the literals true in all its models; applied
to $U_i$ they give unit clauses, and applied to the cofactor of $U_i$ by a literal $\ell$ they give
every binary clause with premise $\ell$ in one call. Extraction therefore needs two cofactors and
essential-variable computations per label in the support of the factor.
\opt{label\_implications} selects \opt{UNCONDITIONAL} (from $U_i$), \opt{PER\_SOURCE} (from the union
over targets of each source state, conditioned on $\prev{s}$), or \opt{BOTH}. Since the empty step is
always allowed, there are no positive units, and almost all unconditional clauses are label mutexes
$\neg\ell_1 \vee \neg\ell_2$.

The number of clauses can be quadratic in the number of labels of a factor. \opt{label\_implications\_max}
stops the extraction at a number of clauses (also within a factor) and
\opt{label\_implications\_time\_limit} after a number of seconds; the clauses found so far are kept,
since each is sound on its own. \opt{label\_implications\_bicliques} encodes the binary clauses by a
greedy biclique cover \parencite{rintanen-ecai2006}: a biclique $A \times B$ of conflicting literals
becomes $a \rightarrow z$, $b \rightarrow \neg z$ with one new variable per step.

\subsection{State clauses}

With $U_i(s) = \bigvee_{s'} \B_i(s,s')$ and $V_i(s') = \bigvee_{s} \B_i(s,s')$ (split BDDs only):
\begin{description}
  \item[\opt{source\_units}] $\prev{s} \rightarrow \pm\ell$ for the literals essential in $U_i(s)$;
  \item[\opt{target\_units}] $\nxt{s'} \rightarrow \pm\ell$ for those essential in $V_i(s')$;
  \item[\opt{pair\_units}] $\prev{s} \wedge \nxt{s'} \rightarrow \pm\ell$ for those essential in
        $\B_i(s,s')$ but not already source or target units: the part of domain consistency Tseitin
        misses for the active BDD;
  \item[\opt{cross\_units}] for two factors sharing labels, $\prev{s_1} \wedge \prev{s_2} \rightarrow
        \pm\ell$ for literals essential in $U_{i}(s_1) \wedge U_{j}(s_2)$ but in neither alone, and the
        same for next states; at most \opt{cross\_pair\_limit} state pairs per factor pair, factor
        pairs with most shared labels first. Because the empty step is always allowed, the
        conjunction is never empty: this yields forced labels, never state mutexes.
\end{description}

\paragraph{Leave and enter clauses.} A factor that starts a step in $s$ and ends elsewhere made its
first move with a label that leaves $s$. A label $\ell$ without any transition from $s$ (not even a
self-loop) can only fire after the factor left $s$, i.e., after a leaving label earlier in the label
order. With \opt{leave\_clauses}:
\[
  \prev{s} \wedge \neg\nxt{s} \rightarrow \bigvee_{\ell \in \mathit{leave}(s)} \ell
  \qquad\text{and}\qquad
  \prev{s} \wedge \ell \rightarrow P_s(\ell)
\]
for every label $\ell$ without a transition from $s$, where $P_s(\ell)$ is an auxiliary variable
implying that some leaving label before $\ell$ fires (one prefix chain per state). \opt{enter\_clauses}
is symmetric for the end state: a label that cannot end in $s'$ needs a later entering label.
\opt{move\_clause\_parts} emits only the frame part (first clause) or only the order part (second),
and \opt{move\_clauses\_max} bounds the number of (state, label) entries. The movers use the same
relevance test as the construction: a label with a self-loop everywhere may still move the factor.

\paragraph{Budget.} The state clauses are computed in phases, cheapest and most useful first:
leave/enter clauses, source/target units (with the unions the later phases need), pair units, cross
units. Within the overall budget \opt{state\_clauses\_time\_limit}, every enabled phase still to come
keeps \opt{state\_clauses\_min\_phase} seconds: phase $k$ stops at
$\mathit{budget} - \mathit{min\_phase} \cdot |\{\text{enabled phases after } k\}|$. What a phase
found before it stopped is kept.

\subsection{Goal condition on the last step}

\begin{proposition}
If a task has a plan with at most $m$ time steps, it has one with at most $m$ steps whose last
action moves some factor from a non-goal state into a goal state.
\end{proposition}
\begin{proof}
Let $a$ be the last action of a plan. Every factor is in a goal state after $a$. If $a$ moved no
factor from a non-goal into a goal state, then every factor was already in a goal state before $a$:
it was either not moved by $a$ or moved from a goal state. Removing $a$ yields a shorter plan with
at most as many steps. Repeat until the condition holds.
\end{proof}
With \opt{goal\_last\_label}, a label that cannot move any factor into its goal states anywhere must,
at the last step, be followed by a later label of that step. The clauses use one suffix chain over
the label order ($A_k$: ``a label at position $\geq k$ fires''). The condition holds for every
horizon, also when the last steps are empty, and is therefore compatible with every length strategy.

%% ===========================================================================
\section{Label Orders}\label{sec:orders}

The label order is given by a \opt{LabelOrderFinder}. The simple finders are \opt{label\_order\_linear}
(label ids), \opt{label\_order\_reverse}, \opt{label\_order\_random}, \opt{label\_order\_causal},
\opt{label\_order\_relaxed} (layers of relaxed reachability), \opt{label\_order\_tsort} (a topological
ranking of the labels) and \opt{label\_order\_file}, which reads an order
and orders unmentioned labels with \opt{leftover}.

\paragraph{Plan-optimal orders as a reference.} For a given plan, the order that minimises the number
of steps is a linear ordering problem; the experiment scripts solve it exactly by MaxSAT. Such orders
are tailored to one plan, but show how much the order matters.

\paragraph{Goal chains.} \opt{label\_order\_goal\_chains} approximates a good order without a plan.
For every goal-relevant factor, a breadth-first search in the product of the factor and its causal
ancestors (up to \opt{ancestor\_depth} levels, within \opt{max\_states} and \opt{state\_budget})
yields an abstract plan, i.e., a chain of labels. The chains are merged into one order by a MaxSAT
over all consecutive pairs (exact per strongly connected component up to \opt{exact\_max\_size}, with
\opt{exact\_time\_limit}); components are ordered topologically, ties by chain index and position.
Further options add joint chains for pairs of goals (\opt{goal\_pairs}), insert leftover labels by
relaxed layer (\opt{leftover\_layer}), and control merging and shrinking of the abstractions
(\opt{merge\_strategy}, \opt{shrink}, \opt{time\_budget}), the sources and scope of order pairs
(\opt{pair\_source}, \opt{pair\_scope}), the scheduling of chains (\opt{chain\_edges},
\opt{chain\_schedule}) and a single abstraction over all goals (\opt{global\_abstraction}). The
configuration we use as default is
\begin{center}
\texttt{label\_order\_goal\_chains(ancestor\_depth=2, goal\_pairs=2, leftover\_layer=true,\\
state\_budget=200000, max\_states=20000, exact\_time\_limit=1)}.
\end{center}

%% ===========================================================================
\section{Search Engines and Length Strategies}\label{sec:engines}

\opt{sat} solves one formula per horizon, with horizons from a \opt{length\_strategy}:
\opt{one\_by\_one} ($1,2,3,\dots$), \opt{by\_iteration} (geometric), \opt{by\_increment} (arithmetic),
\opt{length\_list} (explicit) and \opt{constant}. \opt{rintanen} implements Rintanen's algorithm~C
\parencite{rintanen-et-al-aij2006}: several horizons are solved in interleaved threads with a
scheduler, the $i$-th open horizon receiving turns of \opt{scheduler\_interval}$\cdot\gamma^i$.
\opt{incremental\_sat} reuses one solver across horizons and needs assumptions, which kissat
\parencite{kissat} does not provide. Every call reports kissat's conflicts, decisions and
propagations, which are deterministic for a given formula and seed and therefore a noise-free
measure of search effort.

%% ===========================================================================
\section{Empirical Notes}\label{sec:empirical}

The following observations come from the Autoscale benchmark set \parencite{torralba-et-al-icaps2021}
(1260 agile tasks, 300\,s, 2.5\,GB, kissat), restricted where noted to the 303 tasks solved by at
least one configuration and not solved by the transformation alone.
\begin{itemize}
  \item The path encodings remove 92--96\,\% of the conflicts but enlarge the formula five- to
        sixfold; overall they are slower, except in single domains (childsnack).
  \item All state clauses with a phased 10\,s budget reduce conflicts to about 0.3 and total time to
        about 0.8 of the plain Tseitin encoding (geometric means on tasks of at least 5\,s).
  \item Capped label implications are the cheapest robust improvement.
  \item Coverage differences of a few tasks at the time limit are dominated by noise: two identical
        runs of the same configuration differ by 2--3 solved tasks, with identical conflict
        counts and run times varying by a factor of 0.83--1.31 (5th to 95th percentile).
  \item Kissat's bounded variable elimination interacts with redundant clauses. In single domains it
        turns easy UNSAT proofs into very hard ones (rovers), but disabling it slows every
        configuration down by 11--21\,\%.
\end{itemize}

\printbibliography

\appendix

%% ===========================================================================
\section{Options}

The tables below are generated from the C++ sources by \texttt{extract\_options.py}; option names,
defaults and descriptions are those of the code.

\subsection{\texorpdfstring{\bddsat}{bdd\_sat}}
% generated by extract_options.py from src/search/sat/bdd_encoding.cc -- do not edit
\begin{longtable}{@{}p{0.27\linewidth}p{0.13\linewidth}p{0.54\linewidth}@{}}
\toprule
option & default & meaning \\
\midrule
\endhead
\texttt{one\_step\_only} & \texttt{true} & only allow one real transition per factor and time step (plus self loops before and after it) instead of arbitrary label sequences \\
\texttt{combinebdds} & \texttt{false} & use one BDD per factor over (state, next state, labels) instead of one BDD per (factor, source, target) triple conditioned on the state variables \\
\texttt{bdd\_size\_limit} & \texttt{-1} & per-factor budget of BDD nodes. Oversized (source,target) BDDs fall back to the one-step BDD. -1 means no limit. Requires one\_step\_only=false \\
\texttt{impltseitsin} & \texttt{true} & use implicational (one-directional) Tseitin encoding instead of a bi-implicational one \\
\texttt{path\_encoding} & \texttt{false} & translate each BDD with a path-based encoding (NNFPath, Abio et al. CPAIOR 2016) instead of Tseitin: variables for the nodes and edges on the root-to-1 path, plus clauses that make unit propagation domain consistent on the BDD's variables. About three variables per node and no sharing between BDDs; best with combinebdds \\
\texttt{path\_sharing} & \texttt{false} & with path\_encoding and split BDDs: one path encoding per factor over all its (source, target) BDDs, which share nodes; the BDD of the active pair is selected by an auxiliary variable for previous\_state(s) \& next\_state(s'). Without it every pair's BDD gets its own path encoding \\
\texttt{label\_implications\_bicliques} & \texttt{false} & encode the binary label\_implications clauses by a greedy biclique cover (Rintanen, ECAI 2006): a biclique A x B of conflicting literals becomes a -\textgreater{} z, b -\textgreater{} -z with one new variable per step instead of |A|*|B| binary clauses \\
\texttt{combine\_ite} & \texttt{false} & with combinebdds: build the factor BDD as if-then-else chains over the state variables instead of conjoining full one-hot minterms. Equal on one-hot assignments, which the exactly-one constraints guarantee; |S|\^{}2 instead of |S|\^{}3 BDD operations \\
\texttt{omitforcedvariables} & \texttt{true} & do not create Tseitin variables for BDD nodes whose in-degree is at most forcedvariablesthreshold \\
\texttt{forcedvariablesthreshold} & \texttt{100} & in-degree up to which omitforcedvariables applies \\
\texttt{cutbdds} & \texttt{false} & fixpoint 'cutting' of the BDDs across factors for stronger constraints \\
\texttt{cut\_mode} & \texttt{AND} & with cutbdds: how a constraint is applied \emph{(AND, AND\_LIMITED, RESTRICT, GLOBAL\_AND, GLOBAL\_AND\_LIMITED, GLOBAL\_RESTRICT, AND\_CAPPED)} \\
\texttt{cut\_limit\_factor} & \texttt{1} & with cut\_mode=AND\_LIMITED \\
\texttt{report\_label\_conjunction} & \texttt{false} & diagnostics: conjoin the any-transition label BDDs of all factors (every label set permitted by every factor in one step) under label\_conjunction\_limit new nodes per step, and report size, time and how many factors fit \\
\texttt{label\_conjunction\_limit} & \texttt{10000000} & node limit for report\_label\_conjunction \\
\texttt{label\_implications} & \texttt{NONE} & add clauses implied by the factor BDDs, so unit propagation sees dependencies the Tseitin translation does not propagate \emph{(NONE, UNCONDITIONAL, PER\_SOURCE, BOTH)} \\
\texttt{label\_implications\_max} & \texttt{-1} & at most this many implied clauses of each kind (unconditional, per source); -1 = all. Extraction stops as soon as the cap is reached; the clauses found so far are kept \\
\texttt{label\_implications\_time\_limit} & \texttt{-1} & stop extracting implied clauses after this many seconds and keep the ones found so far (all of them are sound); -1 = no limit \\
\texttt{source\_units} & \texttt{false} & split BDDs: prev\_s -\textgreater{} +-l for the labels forced in every label set the factor permits from s \\
\texttt{target\_units} & \texttt{false} & split BDDs: next\_s' -\textgreater{} +-l for the labels forced in every label set leading into s' \\
\texttt{pair\_units} & \texttt{false} & split BDDs: prev\_s \& next\_s' -\textgreater{} +-l for the labels the (s, s') BDD forces (not already source/target units) \\
\texttt{cross\_units} & \texttt{false} & split BDDs: for two factors sharing labels, prev\_s1 \& prev\_s2 -\textgreater{} +-l (and the same for next states) for labels forced only by both together; factor pairs with at most cross\_pair\_limit state pairs, most shared labels first \\
\texttt{cross\_pair\_limit} & \texttt{10000} & cross\_units: state pairs per factor pair \\
\texttt{leave\_clauses} & \texttt{false} & prev\_s \& -next\_s -\textgreater{} some label leaving s fires; prev\_s \& l -\textgreater{} a label leaving s fires before l (in the label order) for every label l without a transition from s \\
\texttt{enter\_clauses} & \texttt{false} & next\_s' \& -prev\_s' -\textgreater{} some label entering s' fires; next\_s' \& l -\textgreater{} a label entering s' fires after l for every label l without a transition into s' \\
\texttt{state\_clauses\_time\_limit} & \texttt{30} & overall budget for the state clauses, computed in phases leave/enter, source/target units, pair units, cross units; what a phase found before it stopped is kept; -1 = no limit \\
\texttt{state\_clauses\_min\_phase} & \texttt{2} & seconds of the budget reserved for each enabled phase still to come: phase k stops at budget - min\_phase * (enabled phases after k) \\
\texttt{move\_clause\_parts} & \texttt{BOTH} & leave/enter clauses: which of the two parts to emit \emph{(BOTH, FRAME, ORDER)} \\
\texttt{move\_clauses\_max} & \texttt{1000000} & leave/enter clauses: at most this many (state, blocked label) entries; -1 = all \\
\texttt{goal\_last\_label} & \texttt{false} & at the last step, a label that cannot move any factor into its goal states must be followed by a later label (a plan whose last action achieves no goal fact can be shortened) \\
\texttt{cut\_min\_overlap} & \texttt{0} & cutbdds (pairwise modes): only propagate from S to T if the labels they share make up at least this fraction of T's labels \\
\texttt{cube\_encoding\_max} & \texttt{-1} & encode a BDD with at most this many paths to 0 directly, one clause per path (cube of its negation) together with the conditions, without Tseitin variables; larger BDDs keep the Tseitin encoding. -1 = off (Tseitin only) \\
\texttt{coverbdds} & \texttt{false} & search for always-true/false state -\textgreater{} +-label implications and report them. Diagnostic only: this does not change the generated formula \\
\texttt{force\_at\_least\_one\_action} & \texttt{false} & force that every time step contains at least one action \\
\texttt{alternate\_label\_orders} & \texttt{false} & use two label orderings instead of one: the order from label\_order for the odd-numbered time steps and its reverse for the even ones. The label order is the BDD variable order, so this means building and encoding two sets of BDDs, and a step's label variables are then constrained in the order belonging to that step \\
\texttt{report\_label\_implications} & \texttt{false} & mine each factor's label BDD for dependencies that hold in every legal label set (l, -l, l -\textgreater{} l', l -\textgreater{} -l') and report how many there are. Diagnostic only: it does not change the formula \\
\texttt{report\_factor\_statistics} & \texttt{false} & emit one BDDSTAT line per factor with its sizes and construction time. Off by default: it is one line per factor and there can be hundreds. The single construction\_ok/construction\_failed summary is always emitted \\
\texttt{cudd\_init\_nodes} & \texttt{16000000} & initial size of the Cudd unique table, divided by the number of labels \\
\texttt{cudd\_cache\_size} & \texttt{16000000} & initial number of Cudd cache slots. Allocated up front and independent of task size, so it dominates memory on small tasks \\
\texttt{cudd\_max\_memory\_mb} & \texttt{0} & hard cap on Cudd's memory in MB. 0 lets Cudd decide from the machine's RAM, which ignores any per-run budget \\
\texttt{bdd\_init\_time\_limit} & \texttt{-1} & seconds allowed for building the BDDs. When it is used up the run stops with 'BDDSTAT construction\_failed ... reason time\_limit' instead of grinding on until the driver kills it. -1 means no limit \\
\texttt{bdd\_node\_limit} & \texttt{-1} & how many live Cudd nodes the construction may use, reported the same way. -1 means no limit \\
\texttt{label\_order} & \texttt{label\_order\_linear()} & order in which labels may be applied within one time step. Unless bdd\_variable\_order is given, this is also the BDD variable order \\
\texttt{bdd\_variable\_order} & \texttt{} & order of the label variables in the BDDs, independently of label\_order. It changes only the size of the BDDs, not which label sets they allow. Omitted: the label order \\
\texttt{bdd\_reorder} & \texttt{NONE} & Cudd variable reordering, applied once after the BDDs are built (and after cutbdds), before they are Tseitin-translated. The labels keep their meaning; only the variable levels change \emph{(NONE, SIFT, SIFT\_CONVERGE, SYMM\_SIFT, GROUP\_SIFT)} \\
\texttt{bdd\_reorder\_dynamic} & \texttt{false} & also let Cudd reorder automatically during construction, with the bdd\_reorder method \\
\texttt{dump\_bdd\_variable\_order} & \texttt{none} & write the final BDD variable order of the labels (after reordering, top level first) to this file, one label id per line, in the format label\_order\_file reads; 'none' = do not. The option parser lower-cases the path. \\
\texttt{dump\_label\_order} & \texttt{none} & write the (primary) label order to this file, one label id per line; 'none' = do not. The option parser lower-cases the path. \\
\bottomrule
\end{longtable}


\subsection{\texorpdfstring{\opt{label\_order\_goal\_chains}}{label\_order\_goal\_chains}}
% generated by extract_options.py from src/search/task_utils/label_order_goal_chains.cc -- do not edit
\begin{longtable}{@{}p{0.27\linewidth}p{0.13\linewidth}p{0.54\linewidth}@{}}
\toprule
option & default & meaning \\
\midrule
\endhead
\texttt{max\_states} & \texttt{50000} & reachable-state cap per product; ancestors are dropped weakest first above it \\
\texttt{exact\_max\_size} & \texttt{150} & components with at most this many labels are ordered by exact MaxSAT \\
\texttt{exact\_conflicts} & \texttt{200000} & kissat conflict limit per MaxSAT call \\
\texttt{exact\_time\_limit} & \texttt{10} & seconds per component for the MaxSAT \\
\texttt{leftover} & \texttt{label\_order\_relaxed()} & order for the labels no chain mentions \\
\texttt{verbose} & \texttt{false} & report per-goal-factor details \\
\texttt{ancestor\_depth} & \texttt{1} & causal-graph distance up to which ancestors join a goal factor's product (nearest and strongest first; the state cap drops the rest) \\
\texttt{subgoal\_depth} & \texttt{0} & recursion depth of subgoal chains: a chain label with a precondition on a factor outside its product gets a chain for reaching that precondition, placed before it; 0 = off \\
\texttt{plans\_per\_goal} & \texttt{1} & candidate abstract plans per chain; \textgreater{}1 (or plan\_slack\textgreater{}0) enumerates plans and selects one per chain jointly \\
\texttt{plan\_slack} & \texttt{0} & candidate plans may be this much longer than the shortest \\
\texttt{all\_plans} & \texttt{false} & with plans\_per\_goal/plan\_slack: merge all candidate plans as separate chains (each votes with its consecutive pairs) instead of selecting one per chain \\
\texttt{selection\_rounds} & \texttt{5} & rounds of choosing each chain's best candidate under the current order and re-merging \\
\texttt{leftover\_layer} & \texttt{false} & insert leftover labels between the chain labels by relaxed-reachability layer instead of appending them \\
\texttt{state\_budget} & \texttt{-1} & product states the searches beyond the default may explore: goal chains with ancestor\_depth \textgreater{} 1, subgoal chains and pair chains, in that order of generation. Once used up, goal chains fall back to the default depth-1 computation (which is never budgeted) and no further subgoal or pair chains are computed; -1 = unlimited \\
\texttt{work\_budget} & \texttt{-1} & like state\_budget, but counted in expanded states plus candidate labels tested, which tracks the actual cost where states have many applicable labels; -1 = unlimited \\
\texttt{deep\_giveup} & \texttt{0} & with ancestor\_depth \textgreater{} 1: stop trying deep ancestors for the remaining goals after this many consecutive goals fell back to depth 1 (where deep products are too big, every attempt is wasted budget); 0 = never \\
\texttt{goal\_pairs} & \texttt{0} & also plan jointly for each goal factor and up to this many other goal factors sharing most ancestors with it, adding those chains; 0 = off \\
\texttt{merge\_strategy} & \texttt{NONE} & how a goal factor's product is built. Anything but NONE merges one factor at a time and prunes unreachable and dead-end states after every merge (exact); a merge exceeding max\_states stops merging for that goal \emph{(NONE, RANKED, CEGAR, DFP\_ROOT, DFP)} \\
\texttt{merge\_other\_goals} & \texttt{false} & with merge\_strategy: the product's goal also requires every merged factor's own goal (as NONE does); default: the root's goal only \\
\texttt{merge\_until\_covered} & \texttt{false} & with merge\_strategy=RANKED: after the ancestors up to ancestor\_depth, also merge deeper ancestors that no earlier goal's product contains yet \\
\texttt{merge\_skip\_too\_big} & \texttt{false} & with merge\_strategy RANKED or CEGAR: a merge exceeding max\_states skips that factor and goes on with the next one instead of ending the goal's merging (DFP always skips) \\
\texttt{merge\_build\_factor} & \texttt{1} & with merge\_strategy: max\_states applies to the product after pruning; its construction may reach merge\_build\_factor * max\_states states \\
\texttt{shrink} & \texttt{NONE} & with merge\_strategy: how products are shrunk (Variant D) \emph{(NONE, EXACT, PRUNE, PAIRS, GREEDY)} \\
\texttt{prune\_slack} & \texttt{0} & with shrink=PRUNE: how much longer than the shortest plan kept plans may be \\
\texttt{time\_budget} & \texttt{-1} & seconds for the whole order computation; once used up, remaining goals get the cheap depth-1 chain, and no more merges, subgoal or pair chains are computed; -1 = unlimited \\
\texttt{pair\_source} & \texttt{PLAN} & which consecutive label pairs of the goal products the MaxSAT sees; needs a merge\_strategy \emph{(PLAN, OPTIMAL, SLACK, ALL)} \\
\texttt{pair\_slack} & \texttt{1} & with pair\_source=SLACK \\
\texttt{pair\_scope} & \texttt{ALL} & with pair\_source: where the state-induced pairs are used \emph{(ALL, LEFTOVER)} \\
\texttt{chain\_tiebreak} & \texttt{CHAIN} & order of independent components (no pair between them) in the merged order \emph{(CHAIN, POSITION)} \\
\texttt{chain\_edges} & \texttt{BASE} & edges of the precedence graph before the SCCs are formed (option C): extra edges can join chains into larger components that the MaxSAT orders \emph{(BASE, DEPENDENCE, INTERACTION)} \\
\texttt{dep\_weight} & \texttt{1} & weight of one DEPENDENCE / INTERACTION edge \\
\texttt{chain\_schedule} & \texttt{CHAIN} & how components without an edge between them are ordered \emph{(CHAIN, INTERACTION, SIMULATE)} \\
\texttt{global\_abstraction} & \texttt{NONE} & one abstraction over all goals: all goal factors and their transitive ancestors are merged, the goal is the conjunction of all goals, and the product is shrunk with global\_shrink beyond global\_max\_states \emph{(NONE, ONLY, ADD)} \\
\texttt{global\_shrink} & \texttt{GREEDY} & shrinking of the global product (see shrink) \emph{(NONE, EXACT, PRUNE, PAIRS, GREEDY)} \\
\texttt{global\_max\_states} & \texttt{0} & size limit of the global product; 0 = max\_states \\
\texttt{global\_samples} & \texttt{1} & abstract plans taken from the global product: the first is the usual shortest one, the others shortest plans with random tie-breaking \\
\texttt{report\_flaws} & \texttt{false} & diagnostics: per goal chain, the product's factors and the base factors that refute the abstract plan (factor@position); per label, the factors it moves and has preconditions on \\
\texttt{pair\_uncovered} & \texttt{false} & with pair\_scope=LEFTOVER: also take the pairs of all paths of every factor that no goal product contains \\
\texttt{chain\_weight} & \texttt{1} & weight of one occurrence of a chain pair (state-induced pairs: 1..pair\_weight); large values keep the chains' order and let the state pairs place the rest \\
\texttt{exact\_max\_weight} & \texttt{-1} & components whose pairs weigh more than this in total get only the heuristic order (the MaxSAT encoding grows with the weight); -1 = no limit \\
\texttt{pair\_weight} & \texttt{4} & state-induced pairs get weight 1..pair\_weight per product, by the number of states they occur at (a chain pair counts 1 per occurrence) \\
\texttt{max\_merges} & \texttt{-1} & with merge\_strategy: at most this many merges per goal factor; -1 = unlimited \\
\texttt{leftover\_support} & \texttt{false} & place leftover labels depth-first before the chain labels they support (Balyo's topological ranking seeded with the chain order) instead of at the end \\
\bottomrule
\end{longtable}


\subsection{Other label order finders}
\opt{label\_order\_linear}, \opt{label\_order\_reverse}, \opt{label\_order\_random},
\opt{label\_order\_causal} and \opt{label\_order\_relaxed} have no options.
\paragraph{\opt{label\_order\_file}}
% generated by extract_options.py from src/search/task_utils/label_order_finder.cc -- do not edit
\begin{longtable}{@{}p{0.27\linewidth}p{0.13\linewidth}p{0.54\linewidth}@{}}
\toprule
option & default & meaning \\
\midrule
\endhead
\texttt{filename} & \texttt{} & label order, one label id of the transformed task per line; ';' starts a comment. NOTE: the option parser lower-cases the whole --search string, so the path must be lower-case \\
\texttt{leftover} & \texttt{label\_order\_relaxed()} & order for the labels the file does not mention \\
\bottomrule
\end{longtable}

\paragraph{\opt{label\_order\_tsort}}
% generated by extract_options.py from src/search/task_utils/label_order_finder.cc -- do not edit
\begin{longtable}{@{}p{0.27\linewidth}p{0.13\linewidth}p{0.54\linewidth}@{}}
\toprule
option & default & meaning \\
\midrule
\endhead
\texttt{goal\_first} & \texttt{false} & start the depth-first search from the labels that move a goal factor into a goal state \\
\bottomrule
\end{longtable}


\subsection{Search engines}
\paragraph{\opt{sat}}
% generated by extract_options.py from src/search/search_engines/plugin_sat.cc -- do not edit
\begin{longtable}{@{}p{0.27\linewidth}p{0.13\linewidth}p{0.54\linewidth}@{}}
\toprule
option & default & meaning \\
\midrule
\endhead
\texttt{continue\_after\_first\_plan} & \texttt{false} & Normally the planner stops once the first plan is found. If set to true, continue search even if a plan as been found. \\
\texttt{length\_strategy} & \texttt{one\_by\_one()} & strategy to determine plan lengths \\
\texttt{encoder} & \texttt{label\_sat()} & type of formula to use for encoding \\
\texttt{step\_time\_limit} & \texttt{-1} & time limit for each step of the SAT run. Defaults to -1, which means no limit \\
\texttt{step\_time\_limits} & \texttt{[]} & per-call time limits in seconds: the i-th SAT call (i-th length of the strategy) gets the i-th entry, later calls step\_time\_limit. Empty = step\_time\_limit for all \\
\texttt{solver\_quiet} & \texttt{false} & if possible try to put the SAT solver into quiet mode (less output to parse for experiments) \\
\bottomrule
\end{longtable}

\paragraph{\opt{rintanen}}
% generated by extract_options.py from src/search/search_engines/plugin_rintanen.cc -- do not edit
\begin{longtable}{@{}p{0.27\linewidth}p{0.13\linewidth}p{0.54\linewidth}@{}}
\toprule
option & default & meaning \\
\midrule
\endhead
\texttt{length\_strategy} & \texttt{one\_by\_one()} & strategy to determine plan lengths \\
\texttt{encoder} & \texttt{label\_sat()} & type of formula to use for encoding \\
\texttt{max\_parallel\_calls} & \texttt{20} & maximum number of SAT calls run in parallel \\
\texttt{scheduler\_interval} & \texttt{1} & interval for the schedule in seconds \\
\texttt{scheduler\_interval\_ms} & \texttt{-1} & interval for the schedule in milliseconds; overrides scheduler\_interval when \textgreater{}= 0 \\
\texttt{gamma} & \texttt{1.0} & geometric time share: the i-th open length (0 = shortest) gets turns of interval * gamma\^{}i; 1 = round-robin \\
\texttt{trace\_schedule} & \texttt{false} & print one RSCHED line per scheduling event (turn, encoded, sat, unsat) \\
\texttt{memory\_limit\_mb} & \texttt{3500} & memory limit in MBs \\
\texttt{schedule\_formula\_as\_one} & \texttt{true} & schedule the creation of a formula as one item (helps keeping to the memory limit) \\
\texttt{solver\_quiet} & \texttt{false} & if possible try to put the SAT solver into quiet mode (less output to parse for experiments) \\
\bottomrule
\end{longtable}


\subsection{Length strategies}
\paragraph{\opt{one\_by\_one}}
% generated by extract_options.py from src/search/sat/length_strategy.cc -- do not edit
\begin{longtable}{@{}p{0.27\linewidth}p{0.13\linewidth}p{0.54\linewidth}@{}}
\toprule
option & default & meaning \\
\midrule
\endhead
\texttt{maximum\_length} & \texttt{-1} & last length to try; -1 = no limit \\
\bottomrule
\end{longtable}

\paragraph{\opt{by\_iteration}}
% generated by extract_options.py from src/search/sat/length_strategy.cc -- do not edit
\begin{longtable}{@{}p{0.27\linewidth}p{0.13\linewidth}p{0.54\linewidth}@{}}
\toprule
option & default & meaning \\
\midrule
\endhead
\texttt{start\_length} & \texttt{5} & Start value for plan length. Further lengths will be start\_length * (multiplier)\^{}iteration. \\
\texttt{maximum\_iteration} & \texttt{-1} & Number of iterations to be performed. If set to -1, there is no limit.  \\
\texttt{multiplier} & \texttt{1.41} & Multiplier for iteration. \\
\texttt{strictly\_increasing} & \texttt{false} & if the formula rounds to a length not larger than the previous one, use previous+1 instead (small multipliers otherwise repeat lengths) \\
\bottomrule
\end{longtable}

\paragraph{\opt{by\_increment}}
% generated by extract_options.py from src/search/sat/length_strategy.cc -- do not edit
\begin{longtable}{@{}p{0.27\linewidth}p{0.13\linewidth}p{0.54\linewidth}@{}}
\toprule
option & default & meaning \\
\midrule
\endhead
\texttt{start\_length} & \texttt{1} & first length \\
\texttt{increment} & \texttt{1} & added per iteration \\
\texttt{maximum\_iteration} & \texttt{-1} & Number of iterations to be performed. If set to -1, there is no limit. \\
\bottomrule
\end{longtable}

\paragraph{\opt{length\_list}}
% generated by extract_options.py from src/search/sat/length_strategy.cc -- do not edit
\begin{longtable}{@{}p{0.27\linewidth}p{0.13\linewidth}p{0.54\linewidth}@{}}
\toprule
option & default & meaning \\
\midrule
\endhead
\texttt{lengths} & \texttt{} & the lengths to try, in this order \\
\bottomrule
\end{longtable}


\end{document}
